%%%%%%%%%%  scan idx 17  |  folio 4  %%%%%%%%%%
(tranformant la classe d'un cycle premier $Z$ en $\cl(\underline O_Z)$,) dont il s'avère ensuite (Exp.\ XIV 4.2) que le noyau est un sous-groupe de torsion. De plus, on montre dans [RRR] comment les images dans $\Gr^{\bullet}_{\mathrm{top}}(X)$ des classes de Chern $c_i(F)$ peuvent se calculer directement en termes de $\cl(F)\in K(X)$ et des opérations $\lambda^i$ de puissance extérieure dans $K(X)$.

Ceci montre que, pour $X$ schéma noethérien régulier muni d'un Module inversible ample, l'anneau $\Gr^{\bullet}_{\mathrm{top}}(X)$ est un substitut commode pour l'anneau de Chow hypothétique $A(X)$~; il a de plus l'avantage sur ce dernier que c'est directement dans cet anneau que se font certains calculs de [RRR] aboutissant à des formules importantes dans cet anneau, formules qui même pour $X$ projectif et lisse sur un corpa ne sont pas établies à l'heure actuelle dans $A(X)$ lui-même (Exp.\ XIV, 4.3 et 4.4). Pour justifier entièrement ce point de vue, il faudrait cependant résoudre le problème signalé plus baut de la compatibilité de la filtration topologique de $K(X)$ avec sa structure d'anneau, problème qui semble essentiellement lié au «~moving lemma~» de Chow, qui lui-même est l'ingrédient technique essentiel de la théorie de l'anneau de Chow que nous avons prétendu pouvoir éviter.

\subsubsection*{2.3.}
Pour éviter ce problème, on munit $K(X)$ d'une deuxième filtration, plus fine que la filtration topologique, et qui peut se décrire en termes purement algébriques à partir de la structure de $\lambda$-anneau augmenté de $K(X)$. La définition de cette filtration est suggérée par l'expression donnée dans [RRR] des classes de Chern d'un $x\in K(X)$ dans $\Gr^{\bullet}_{\mathrm{top}}(X)$,
\begin{equation}\tag{2.3}
c_i(x)=\gamma^i(x-\epsilon(x))\mod \operatorname{Filt}^{i+1}_{\mathrm{top}}K(X),
\end{equation}
où $\gamma^i$ désigne une certaine combinaison linéaire des $\lambda^j$ ($j=i$) explicitée dans loc. cit., $\epsilon$ étant l'augmentation. On montre que $\gamma^i(x-\epsilon(x))$ est bien de filtration topologique $\ge i$, ce qui donne un sens à la formule (2.3). On définit alors la $\lambda$-\underline{filtration} de $K(X)$ par les idéaux $\operatorname{Filt}^i K(X)$, où $\operatorname{Filt}^i K(X)$ est formé par les sommes d'expression de la forme
\[
\gamma^{i_1}(x_1-\epsilon(x_1))\cdots \gamma^{i_k}(x_k-\epsilon(x_k)),\qquad i_1+\cdots+i_k\ge i.
\]
L'anneau gradué associé à cette filtration sera noté $\Gr^\bullet(X)$, et c'est lui qui remplacera dans le Séminaire l'anneau de Chow de $X$. On peut d'ailleurs montrer que dans le cas envisagé ici l'homomorphisme canonique
\begin{equation}\tag{2.4}
\Gr^\bullet(X)\longrightarrow \Gr^{\bullet}_{\mathrm{top}}(X)
\end{equation}
